\section{Techniken zur Messung von Raum-Impulsantworten}
\label{sec:ir_msrmt}
Um die akustischen Eigenschaften eines physischen Raumes in Form einer Impulsantwort abbilden zu können, wird der Raum
zunächst durch einen Messton angeregt. Dieser muss das gesamte Frequenzspektrum, auf das die IR später angewandt werden soll, enthalten \cite{smpm2001}.
Der im Raum resultierende Raumschall (der aus den Eigenschaften des Raumes sowie dem anregenden Ton entstanden ist) wird aufgenommen.
Das aufgenommene Signal ist das Ergebnis der Transferfunktion des gemessenen Raumes,
angewandt auf den Messton. Da Raumhall-Effekte sich wie in \ref{sec:irl_lti} dargestellt durch ihre Impulsantwort annähern lassen, kann die IR des Raumes
mithilfe entsprechender Verfahren aus diesem Signal berechnet werden \cite{novak2016}. Dabei macht es für die am Ende des Entfaltungsprozesses resultierende Impulsantwort
einen entscheidenden Unterschied, wie das Signal, mit dem der Raum angegeregt wird, beschaffen ist. In der Vergangenheit wurden oft sogenannte
pulsive Schallquellen genutzt \cite{farina2007b}, etwa Schüsse oder platzende Luftballons. Mit der steigenden Verbreitung digitaler Systeme
haben sich Mitte der 1980er Jahre zwei Schallquellen, die Maximum-Length-Sequence (kurz \textit{MLS}) und die Time Delay Spectrometry (kurz \textit{TDS}),
bei der ein linearer Sinus-Sweep über das zu messende Frequenzspektrum zum Einsatz kommt, durchgesetzt. Mit diesen Methoden konnten
verbesserter Rauschabstand und flachere Frequenzgänge der Impulsantworten erzielt werden \cite{farina2007b}. Seit Anfang der 2000er wird vermehrt
an Messverfahren mit exponentiellen/logarithmischen Sinus-Sweeps (kurz \textit{ESS/LSS}) geforscht.
Diese Verfahren ermöglichen verbesserte Abbildung der gemessenen Systeme besonders bei Vorhandensein von Nonlinearen oder Zeit-Varianten Eigenschaften \cite{farina2007b}.

\paragraph{Time Delay Spectrometry (TDS)}
TDS-basierte Messverfahren nutzen einen Sinus-Sweep, bei dem die Frequenz des Messsignals über die Messdauer linear durch den zu messenden
Frequenzbereich verändert wird. TDS-Verfahren sind eher für die Messung kurzer Impulsantworten geeignet, etwa denen von Lautsprecher-Systemen \cite{farina2000, holters2009}. 
\paragraph{Minimum Length Sequence (MLS)}
MLS-basierte Messverfahren nutzen periodische Sequenzen von binären Ziffern, die auch als Pseudo-Random Noise bezeichnet werden.
MLS übertragen nur wenig Energie in kurzen Impulsen. Als Ausgleich werden viele Messungen
vorgenommen und das Ergebnise gemittelt. Unter Anwendung entsprechender Prozesse weist das Ergebnis einer MLS-Messung einen besseren Rauschabstand auf, als TDS-Messungen \cite{rife1989}.
Allerdings sind MLS-basierte Messverfahren sehr anfällig gegenüber Nonlinearitäten im gemessenen System und benötigen somit besondere Sorgfalt beim Einstellen der Messpegel \cite{holters2009}.
\paragraph{Exponential/Logarithmic Sine Sweep (ESS/LSS)}
ESS/LSS-basierte Messverfahren nutzen exponentielle Sinus-Sweeps, bei denen die momentäre Frequenz sich in Abhängigkeit der Zeit exponentiell verändert.
Bei MLS- bzw. TDS-basierten Verfahren auftretende Artefake auftretende Artefakte können durch ein von \cite{farina2000} beschriebenes Entfaltungs-Verfahren einfach von der IR getrennt werden können.
Unter Anwendung dieser Methodik ist sind ESS/LSS-Verfahren die geeignetsten Messverfahren für die Messung der Impulsantworten von Räumen.



